%% kapa: exam=1 round=23 year=2012 period=1 form=A subject=civil_law subject_ko=민법 seq=32 number=32
%% source: Q-Net 감정평가사 기출문제 (cst00309) artlSeq=5052578 file=raw/1008/5052578/2012년 제23회 감정평가사 제1차시험 기출문제(1교시 A형).pdf
%% conversion: kordoc 4.24.51 -> tools/gen1.py (auto); review=proofread-sonnet; flags=none
\begin{question}{32}
주위토지통행권에 관한 설명으로 옳지 않은 것은? (다툼이 있으면 판례에 의함)
\begin{choices}
  \item 주위토지통행권이 발생한 후 그 토지에 접하는 공로가 개설됨으로써 이를 인정할 필요성이 없어진 때에는 주위토지통행권이 소멸한다.
  \item 토지의 용도에 필요한 통로가 있는 경우에는 그 통로를 사용하는 것보다 더 편리하다는 이유만으로 주위토지통행권을 인정할 수는 없다.
  \item 주위토지통행권의 본래적 기능 발휘를 위해서 그 통행에 방해가 되는 축조물은 적법하게 설치되었던 것이라 하더라도 철거되어야 한다.
  \item 특별한 사정이 없는 한 주위토지통행권이 성립한 경우, 통행권자는 통행지소유자의 손해를 보상하여야 한다.
  \item 토지사용권이 없는 명의신탁자에게도 주위토지통행권이 인정된다.
\end{choices}
\end{question}
\answer{5}
